\section{Convergence to Minimizer vs. Channel Collapse: Two Different Phenomena}
\label{sec:convergence-vs-collapse}

This section clarifies the distinction between two different concepts: (1) convergence to a global minimizer, and (2) channel collapse.

\subsection{Convergence to Global Minimizer: "If It Converges, Then to Minimizer"}

\textbf{The theorem states:} Under the assumptions (including \emph{Convergence to Limit Point}), if the mean-field dynamics $W(t)$ converge to some limit point $W^*$ as $t \to \infty$, then $W^*$ is a \textbf{global minimizer} of the population loss.

\textbf{What this means:}
\begin{itemize}
\item \textbf{Convergence assumption}: We assume $W(t) \to W^*$ (in a Wasserstein-like sense)
\item \textbf{Conclusion}: $W^*$ minimizes $\mathscr{L}(W) = \mathbb{E}_Z[\mathcal{L}(Y, \hat{y}(X; W))]$
\item \textbf{Key mechanism}: Universal approximation property + convergence $\to$ global optimality
\end{itemize}

\textbf{The proof strategy:}
\begin{enumerate}
\item Show that at the limit point, the gradient is zero: $\partial_t w_i(t) \to 0$
\item Use universal approximation property (frozen random features maintain full support)
\item Conclude that if gradient is zero and function class is dense, then we're at a global minimizer
\end{enumerate}

\textbf{This is guaranteed} (under the assumptions) - if the solution converges, it converges to a global minimizer.

\subsection{Channel Collapse: A Different Issue}

\textbf{Channel collapse} is a separate phenomenon: all $r$ channels learn the \textbf{same function}, i.e., $f_1(t, \cdot) = f_2(t, \cdot) = \cdots = f_r(t, \cdot)$.

\textbf{What this means:}
\begin{itemize}
\item The network still converges to a global minimizer (if convergence assumption holds)
\item But the benefit of having $r$ channels is lost (they're all identical)
\item The low-rank structure becomes redundant
\end{itemize}

\textbf{This is NOT guaranteed} - channels may or may not collapse, depending on the mixing matrix $L$ and initialization.

\subsection{What Causes Channel Collapse?}

\textbf{Channel collapse occurs when all channels receive the same gradient signal:}

\begin{enumerate}
\item \textbf{Isotropic Gaussian mixing}: If $L_{C_2} \sim \mathcal{N}(0, I_r)$ (isotropic), then:
   \[
   (B_1, \ldots, B_r) = c(t) \cdot (f_1, \ldots, f_r) \quad \text{(colinear)}
   \]
   All channels receive \textbf{proportional} gradients, so if they start identically, they stay identical.

\item \textbf{Symmetric initialization}: If all channels start with the same initialization:
   \[
   w_1^0(C_1, 1) = w_1^0(C_1, 2) = \cdots = w_1^0(C_1, r)
   \]
   and the mixing matrix is symmetric/isotropic, then all channels evolve identically.

\item \textbf{Uniform mixing matrix}: If $L$ has no structure (e.g., all entries equal), then all channels receive similar signals and may not differentiate.

\item \textbf{No ReLU gating differentiation}: If the ReLU gate $\mathbf{1}\{H_2 > 0\}$ activates the same neurons for all channels, then $B_k$ become similar for all $k$.
\end{enumerate}

\subsection{What Prevents Channel Collapse?}

\textbf{Mechanisms that enable specialization:}

\begin{enumerate}
\item \textbf{Non-isotropic mixing matrix}: Structured $L$ (e.g., deterministic grid) creates different effective gradients for each channel.

\item \textbf{Asymmetric initialization}: Different initializations for different channels can be amplified by the log-ratio growth mechanism.

\item \textbf{ReLU gating creates competition}: The gate $\mathbf{1}\{H_2 > 0\}$ depends on the sum over channels, creating competitive dynamics where different channels can dominate at different locations.

\item \textbf{Log-ratio growth preserves dominance}: Once a channel establishes dominance, the log-ratio mechanism preserves and amplifies it, forcing other channels to specialize elsewhere.
\end{enumerate}

\subsection{The Key Distinction}

\textbf{Convergence to minimizer:}
\begin{itemize}
\item \textbf{Guaranteed} (under convergence assumption)
\item \textbf{Independent} of channel collapse
\item \textbf{Means}: The network learns the target function optimally
\end{itemize}

\textbf{Channel collapse:}
\begin{itemize}
\item \textbf{Not guaranteed} - depends on $L$ and initialization
\item \textbf{Separate} from convergence to minimizer
\item \textbf{Means}: All channels learn the same function (redundancy)
\end{itemize}

\textbf{Important:} Even if channels collapse, the network can still converge to a global minimizer! The collapse just means we're not using the full capacity of having $r$ channels.

\subsection{Example Scenarios}

\textbf{Scenario 1: Convergence with specialization}
\begin{itemize}
\item $W(t) \to W^*$ (converges)
\item $W^*$ is a global minimizer (guaranteed by theorem)
\item $f_1^*, f_2^*, \ldots, f_r^*$ are different (channels specialized)
\item \textbf{Best case}: Network uses full capacity, channels capture different frequencies/regions
\end{itemize}

\textbf{Scenario 2: Convergence with collapse}
\begin{itemize}
\item $W(t) \to W^*$ (converges)
\item $W^*$ is a global minimizer (still guaranteed by theorem)
\item $f_1^* = f_2^* = \cdots = f_r^*$ (all channels identical)
\item \textbf{Still works}: Network learns the target, but $r$ channels are redundant
\end{itemize}

\textbf{Scenario 3: No convergence}
\begin{itemize}
\item $W(t)$ does not converge (oscillates, diverges, etc.)
\item Theorem doesn't apply (convergence assumption violated)
\item \textbf{Problem}: Can't guarantee global optimality
\end{itemize}

\subsection{What the Theory Guarantees vs. What It Doesn't}

\textbf{Guaranteed:}
\begin{itemize}
\item \textbf{If} $W(t) \to W^*$ (convergence assumption), \textbf{then} $W^*$ is a global minimizer
\item Well-posedness of the dynamics
\item Quantitative approximation bounds
\end{itemize}

\textbf{Not guaranteed:}
\begin{itemize}
\item \textbf{That} $W(t)$ converges (this is an assumption)
\item \textbf{That} channels specialize (depends on $L$ and initialization)
\item \textbf{That} channels don't collapse (isotropic mixing can cause collapse)
\end{itemize}

\subsection{Conclusion}

\textbf{To answer your question:}

\begin{enumerate}
\item \textbf{"If solution converges, then to minimizer"}: Yes, this is guaranteed by the theorem. If $W(t) \to W^*$, then $W^*$ minimizes the population loss (under the assumptions).

\item \textbf{"Collapse can occur due to what?"}: Channel collapse occurs when:
   \begin{itemize}
   \item \textbf{Isotropic Gaussian mixing}: Makes $B_k$ colinear with $f_k$, preventing differentiation
   \item \textbf{Symmetric initialization}: All channels start identically
   \item \textbf{Uniform mixing}: $L$ has no structure to create different gradients
   \item \textbf{No ReLU differentiation}: Gate activates same neurons for all channels
   \end{itemize}

\item \textbf{Key point}: Convergence to minimizer and channel collapse are \textbf{independent}. You can have:
   \begin{itemize}
   \item Convergence with specialization (best case)
   \item Convergence with collapse (still works, but redundant)
   \item No convergence (theorem doesn't apply)
   \end{itemize}
\end{enumerate}

\textbf{The theory guarantees global optimality IF convergence occurs, but does NOT guarantee that channels won't collapse. Channel specialization is enabled by the low-rank structure with non-isotropic mixing, but it's not guaranteed.}
